% Options for packages loaded elsewhere
\PassOptionsToPackage{unicode}{hyperref}
\PassOptionsToPackage{hyphens}{url}
\documentclass[
]{article}
\usepackage{xcolor}
\usepackage{amsmath,amssymb}
\setcounter{secnumdepth}{-\maxdimen} % remove section numbering
\usepackage{iftex}
\ifPDFTeX
  \usepackage[T1]{fontenc}
  \usepackage[utf8]{inputenc}
  \usepackage{textcomp} % provide euro and other symbols
\else % if luatex or xetex
  \usepackage{unicode-math} % this also loads fontspec
  \defaultfontfeatures{Scale=MatchLowercase}
  \defaultfontfeatures[\rmfamily]{Ligatures=TeX,Scale=1}
\fi
\usepackage{lmodern}
\ifPDFTeX\else
  % xetex/luatex font selection
\fi
% Use upquote if available, for straight quotes in verbatim environments
\IfFileExists{upquote.sty}{\usepackage{upquote}}{}
\IfFileExists{microtype.sty}{% use microtype if available
  \usepackage[]{microtype}
  \UseMicrotypeSet[protrusion]{basicmath} % disable protrusion for tt fonts
}{}
\makeatletter
\@ifundefined{KOMAClassName}{% if non-KOMA class
  \IfFileExists{parskip.sty}{%
    \usepackage{parskip}
  }{% else
    \setlength{\parindent}{0pt}
    \setlength{\parskip}{6pt plus 2pt minus 1pt}}
}{% if KOMA class
  \KOMAoptions{parskip=half}}
\makeatother
\usepackage{longtable,booktabs,array}
\usepackage{caption}
\captionsetup[table]{skip=6pt}
\newcounter{none} % for unnumbered tables
\usepackage{calc} % for calculating minipage widths
% Correct order of tables after \paragraph or \subparagraph
\usepackage{etoolbox}
\makeatletter
\patchcmd\longtable{\par}{\if@noskipsec\mbox{}\fi\par}{}{}
\makeatother
% Allow footnotes in longtable head/foot
\IfFileExists{footnotehyper.sty}{\usepackage{footnotehyper}}{\usepackage{footnote}}
\makesavenoteenv{longtable}
\usepackage{graphicx}
\makeatletter
\newsavebox\pandoc@box
\newcommand*\pandocbounded[1]{% scales image to fit in text height/width
  \sbox\pandoc@box{#1}%
  \Gscale@div\@tempa{\textheight}{\dimexpr\ht\pandoc@box+\dp\pandoc@box\relax}%
  \Gscale@div\@tempb{\linewidth}{\wd\pandoc@box}%
  \ifdim\@tempb\p@<\@tempa\p@\let\@tempa\@tempb\fi% select the smaller of both
  \ifdim\@tempa\p@<\p@\scalebox{\@tempa}{\usebox\pandoc@box}%
  \else\usebox{\pandoc@box}%
  \fi%
}
% Set default figure placement to htbp
\def\fps@figure{htbp}
\makeatother
\setlength{\emergencystretch}{3em} % prevent overfull lines
\providecommand{\tightlist}{%
  \setlength{\itemsep}{0pt}\setlength{\parskip}{0pt}}
\usepackage{bookmark}
\IfFileExists{xurl.sty}{\usepackage{xurl}}{} % add URL line breaks if available
\urlstyle{same}
% fallback for those not using the hyperref driver hyperxmp:
\makeatletter
\@ifundefined{xmpquote}{\newcommand{\xmpquote}[1]{#1}}{}
\makeatother
\hypersetup{
  hidelinks,
  pdfcreator={LaTeX via pandoc}}

\author{}
\date{}

\begin{document}

\section{Bab 7: Kegagalan Tata Kelola: D3TLH Dalam Sistem
Perizinan}\label{bab-7-kegagalan-tata-kelola-d3tlh-dalam-sistem-perizinan}

\textbf{CELIOS --- Center of Economic and Law Studies}

\emph{Evaluasi instrumen perlindungan ekologis dan implementasinya dalam
sistem perizinan.}

\begin{center}\rule{0.5\linewidth}{0.5pt}\end{center}

\subsection{Ringkasan Tata Kelola}\label{ringkasan-tata-kelola}

{\def\LTcaptype{none} % do not increment counter
\begin{longtable}[]{@{}
  >{\raggedright\arraybackslash}p{(\linewidth - 4\tabcolsep) * \real{0.3333}}
  >{\raggedright\arraybackslash}p{(\linewidth - 4\tabcolsep) * \real{0.3333}}
  >{\raggedright\arraybackslash}p{(\linewidth - 4\tabcolsep) * \real{0.3333}}@{}}
\toprule\noalign{}
\begin{minipage}[b]{\linewidth}\raggedright
Indikator
\end{minipage} & \begin{minipage}[b]{\linewidth}\raggedright
Kondisi Aktual
\end{minipage} & \begin{minipage}[b]{\linewidth}\raggedright
Deskripsi
\end{minipage} \\
\midrule\noalign{}
\endhead
\bottomrule\noalign{}
\endlastfoot
\textbf{Fungsi Pembatas Daya Dukung} & \textbf{Penerbitan Izin Lanjut} &
Data menunjukkan perlunya penguatan fungsi D3TLH sebagai instrumen
pengaman. Penerbitan izin baru masih berlangsung di kawasan yang
tercatat mengalami tekanan lingkungan tinggi. \\
\textbf{Pengawasan \& Penegakan Hukum} & \textbf{Tantangan Sanksi} &
Evaluasi menunjukkan tantangan dalam penegakan sanksi administratif dan
pengawasan perizinan bagi entitas yang beroperasi tidak sesuai
ketentuan. \\
\end{longtable}
}

\begin{center}\rule{0.5\linewidth}{0.5pt}\end{center}

\subsection{Metodologi Pendekatan \& Pertanyaan
Kritis}\label{metodologi-pendekatan-pertanyaan-kritis}

\textbf{Membaca Hubungan Antara Hasil D3TLH dan Keputusan Perizinan
Aktual}

Analisis di halaman ini menggunakan \textbf{Matriks Analisis Crosstab}
untuk menjawab 3 pertanyaan fundamental: 1. \textbf{Apakah D3TLH
digunakan sebagai dasar keputusan?} 2. \textbf{Apakah D3TLH bersifat
mengikat atau hanya rekomendasi?} 3. \textbf{Apakah D3TLH dapat
diabaikan secara prosedural?}

\emph{Kerangka Pengujian:} Menyilangkan Data Fase Status Ekologis (Aman
/ Tertekan / Kritis) dengan Data Empiris Keputusan Izin yang benar-benar
diterbitkan negara.

\begin{center}\rule{0.5\linewidth}{0.5pt}\end{center}

Dokumen tata ruang dan Daya Dukung Daya Tampung Lingkungan Hidup (D3TLH)
dirancang sebagai instrumen pengendalian investasi agar tidak melampaui
kapasitas daya dukung lingkungan. Penelusuran \emph{timeline} penerbitan
izin di Sulawesi mengindikasikan perlunya penguatan efektivitas
instrumen ini dalam proses perizinan usaha pertambangan dan
infrastruktur pendukungnya.

\begin{center}\rule{0.5\linewidth}{0.5pt}\end{center}

\subsection{7.1 Pembuktian Empiris: Status Ekologis vs Penerbitan
Izin}\label{pembuktian-empiris-status-ekologis-vs-penerbitan-izin}

\textbf{Metode: Spatial Overlay \& Crosstabulation (ESDM x GFW)}

\subsubsection{Metodologi: Evaluasi Kepatuhan D3TLH Berdasarkan Data
Historis}\label{metodologi-evaluasi-kepatuhan-d3tlh-berdasarkan-data-historis}

\textbf{Metode Analisis:} Sub-bab ini menggunakan agregasi berbasis
aturan (\emph{Rule-based Categorization}) untuk membedah ketidaksesuaian
antara status kerusakan lingkungan dengan keputusan administratif
perizinan.

\begin{enumerate}
\def\labelenumi{\arabic{enumi}.}
\tightlist
\item
  \textbf{Model Evaluasi Pelanggaran (\emph{Compliance Modeling}):}

  \begin{itemize}
  \tightlist
  \item
    \textbf{Kategorisasi Status (Binning):} Nilai kerusakan lingkungan
    absolut dibagi ke dalam tiga kelas menggunakan distribusi
    \emph{percentile}: Aman (≤33\%), Tertekan (33-66\%), dan Kritis
    (\textgreater66\%).
  \item
    \textbf{Kuantifikasi Pelanggaran D3TLH:} Mengidentifikasi secara
    kuantitatif apakah pemerintah tetap mengobral Izin Usaha
    Pertambangan (IUP) baru pada wilayah-wilayah yang secara empiris
    terbukti telah berada di fase `Kritis'.
  \end{itemize}
\item
  \textbf{Kalkulasi/Formula Pengolahan:}

  \begin{itemize}
  \tightlist
  \item
    \texttt{Ambang\_Kritis\ =\ Percentile(Deforestasi,\ 0.66)}
  \item
    \texttt{Total\_Izin\_Ilegal\_Ekologis\ =\ SUM(IUP\_Baru)\ WHERE\ Status\_D3TLH\ =\ \textquotesingle{}Kritis\textquotesingle{}}
  \end{itemize}
\item
  \textbf{Variabel \& Fitur Data:}

  \begin{itemize}
  \tightlist
  \item
    \textbf{Variabel Konteks Lingkungan:}
    \texttt{Total\_Deforestasi\_Ha} atau
    \texttt{Deforestasi\_Driver\_Komoditas...} (sebagai basis status
    wilayah)
  \item
    \textbf{Variabel Keputusan Aktor:} \texttt{Jumlah\_Izin\_Baru} dan
    \texttt{Total\_Luas\_Konsesi\_Baru\_Ha}
  \end{itemize}
\item
  \textbf{Dataset \& File:}

  \begin{itemize}
  \tightlist
  \item
    \texttt{data/processed/sulawesi\_izin\_baru\_per\_tahun.csv} dan
    \texttt{sulawesi\_gfw\_master\_1\_dekade\_2014\_2023.csv}
  \end{itemize}
\end{enumerate}

\begin{center}\rule{0.5\linewidth}{0.5pt}\end{center}

Daya Tampung dan Daya Dukung Lingkungan Hidup (D3TLH) dirancang sebagai
instrumen pencegahan dan pengatur batas pengaman ekologis
(\emph{ecological safeguard}). Secara umum, penerbitan izin baru
sepatutnya mempertimbangkan indikator daya dukung lingkungan guna
mengantisipasi degradasi ekosistem.

Penyandingan data deforestasi tahunan dari \emph{Global Forest Watch}
(GFW) dan data perizinan pertambangan dari \emph{Minerba One Data
Indonesia} (MODI) Kementerian ESDM menunjukkan bahwa penerbitan izin
usaha pertambangan baru tetap tercatat pada kurun waktu ketika perubahan
tutupan hutan meningkat. Hal ini terlihat pada tren di wilayah Sulawesi
Tengah dan Tenggara periode 2014-2023.

Kondisi ini menggarisbawahi pentingnya penguatan fungsi dokumen AMDAL,
D3TLH, dan KLHS agar menjadi pertimbangan utama yang mengikat dalam
pengambilan keputusan perizinan, demi menjaga keberlanjutan lingkungan
dan kehidupan masyarakat sekitar.

Matriks Statistik di bawah ini menyajikan perbandingan indikator status
ekologis dan penerbitan izin.

\subsubsection{Konklusi Analisis Kepatuhan D3TLH Berdasarkan Data
Historis}\label{konklusi-analisis-kepatuhan-d3tlh-berdasarkan-data-historis}

Tabel pembuktian di bawah ini mengukur akumulasi penerbitan izin pada
rentang waktu ketika wilayah berstatus Aman, Tertekan, hingga Kritis
secara aktual. Apabila pada status Kritis izin masih diterbitkan, hal
tersebut secara matematis mendiskualifikasi D3TLH sebagai instrumen
perlindungan lingkungan.

{\def\LTcaptype{none} % do not increment counter
\begin{longtable}[]{@{}
  >{\raggedright\arraybackslash}p{(\linewidth - 8\tabcolsep) * \real{0.2000}}
  >{\raggedright\arraybackslash}p{(\linewidth - 8\tabcolsep) * \real{0.2000}}
  >{\raggedright\arraybackslash}p{(\linewidth - 8\tabcolsep) * \real{0.2000}}
  >{\raggedright\arraybackslash}p{(\linewidth - 8\tabcolsep) * \real{0.2000}}
  >{\raggedright\arraybackslash}p{(\linewidth - 8\tabcolsep) * \real{0.2000}}@{}}
\toprule\noalign{}
\begin{minipage}[b]{\linewidth}\raggedright
Status Daya Dukung
\end{minipage} & \begin{minipage}[b]{\linewidth}\raggedright
Kondisi Kerusakan Hutan
\end{minipage} & \begin{minipage}[b]{\linewidth}\raggedright
Seharusnya (Menurut Aturan)
\end{minipage} & \begin{minipage}[b]{\linewidth}\raggedright
Kenyataan di Lapangan
\end{minipage} & \begin{minipage}[b]{\linewidth}\raggedright
Kesimpulan Tata Kelola
\end{minipage} \\
\midrule\noalign{}
\endhead
\bottomrule\noalign{}
\endlastfoot
\textbf{Aman} & Ringan (Hilang 2,649 - 14,542 Ha) & Wajar diterbitkan
izin & 39 Izin Baru Keluar & Normal (Sesuai Aturan) \\
\textbf{Tertekan} & Sedang (Hilang 14,568 - 36,808 Ha) & Izin mulai
direm/dibatasi & 81 Izin Baru Keluar & Anomali (Lampu Kuning) \\
\textbf{Kritis} & Sangat Parah (Hilang 39,633 - 154,169 Ha) &
\textbf{Moratorium / Izin Dilarang!} & \textbf{260 Izin Baru Keluar}
(Termasuk luasan 447,683 Ha) & \textbf{BUKTI PELANGGARAN FATAL} \\
\end{longtable}
}

\begin{quote}
\textbf{Temuan Target:} - \textbf{Fungsi pembatas D3TLH perlu
ditingkatkan} (Terdapat \textbf{260} Izin Baru yang terbit pada periode
berstatus deforestasi tinggi). - \textbf{Diperlukan penguatan integrasi
data lingkungan dalam keputusan perizinan}.
\end{quote}

\textbf{Metodologi Pembuktian (100\% Data-Driven)}

Daftar di bawah ini adalah \textbf{Tabel Irisan (Intersection)} yang
menyatukan Data Satelit (GFW) dan Data Perizinan (ESDM Minerba One).
Sistem secara otomatis melacak dan menarik nama-nama perusahaan yang SK
IUP-nya ditandatangani persis pada Tahun dan Provinsi yang sedang
berstatus Kritis akibat deforestasi.

\paragraph{Tabel Irisan: Daftar Izin IUP Baru yang Diterbitkan di Tengah
Situasi
Kritis}\label{tabel-irisan-daftar-izin-iup-baru-yang-diterbitkan-di-tengah-situasi-kritis}

\emph{(Tabel dapat dilihat pada versi interaktif web).}

\begin{center}\rule{0.5\linewidth}{0.5pt}\end{center}

\subsection{7.2 Tabrakan Hukum: Impunitas dan Pembiaran Operasi
Ilegal}\label{tabrakan-hukum-impunitas-dan-pembiaran-operasi-ilegal}

\textbf{Impunitas Korporasi dan Pembiaran Konflik Struktural di Sektor
Ekstraktif}

\textbf{Metode: Thematic Coding \& Analisis Kasus (LSM / KPA / Tanah
Kita)}

\subsubsection{Metodologi: Pemetaan Impunitas
Korporasi}\label{metodologi-pemetaan-impunitas-korporasi}

\textbf{Metode Analisis:} Sub-bab ini menggunakan agregasi pelaporan
berbasis insiden (\emph{Incident-based Reporting Aggregation}) untuk
mengukur tingkat pembiaran penegakan hukum (impunitas).

\begin{enumerate}
\def\labelenumi{\arabic{enumi}.}
\tightlist
\item
  \textbf{Model Penelusuran Anomali Hukum:}

  \begin{itemize}
  \tightlist
  \item
    \textbf{Kasus Rekam Jejak:} Menyaring dan mengklasifikasikan
    database konflik sengketa lahan, pelanggaran HAM, dan kasus operasi
    ilegal tanpa izin di level tapak.
  \item
    \textbf{Pemetaan Pembiaran (\emph{State Omission}):} Menghitung
    total volume agregat di mana korporasi yang terbukti bermasalah
    secara hukum tetap dipertahankan keberadaan operasinya oleh aparatur
    negara.
  \end{itemize}
\item
  \textbf{Kalkulasi/Formula Pengolahan:}

  \begin{itemize}
  \tightlist
  \item
    \texttt{Total\_Kasus\_Impunitas\ =\ COUNT(Judul\_Kasus)}
  \item
    \texttt{Volume\_Pembiaran\_Sektoral\ =\ SUM(Kasus)\ GROUP\ BY\ Sektor}
  \end{itemize}
\item
  \textbf{Variabel \& Fitur Data:}

  \begin{itemize}
  \tightlist
  \item
    \textbf{Atribut Laporan:} \texttt{Provinsi}, \texttt{Sektor},
    \texttt{Judul\_Kasus}, \texttt{Deskripsi\_Singkat}
  \end{itemize}
\item
  \textbf{Dataset \& File:}

  \begin{itemize}
  \tightlist
  \item
    \texttt{data/processed/sulawesi\_konflik\_hukum.csv}
  \end{itemize}
\end{enumerate}

\begin{center}\rule{0.5\linewidth}{0.5pt}\end{center}

Konsep Daya Tampung dan Daya Dukung Lingkungan Hidup (D3TLH) mengukur
kapasitas daya tahan ekosistem serta daya dukung sosial masyarakat di
sekitar kawasan industri. Kompilasi laporan masyarakat sipil dan
organisasi terkait mencatat adanya sengketa tanah dan dinamika sosial
dalam ekspansi industri ekstraktif.

Hal ini menunjukkan pentingnya kepatuhan perizinan dan penerapan sanksi
administratif secara konsisten. Pengawasan terhadap batas wilayah
perizinan (HGU/IUP) serta pelaksanaan konsultasi publik (FPIC) menjadi
aspek penting dalam tata kelola pertanahan dan lingkungan.

Penguatan koordinasi antar-instansi serta penyelesaian sengketa tenurial
secara adil menjadi langkah krusial untuk memastikan kepastian hukum dan
perlindungan hak masyarakat di wilayah sekitar industri.

{\def\LTcaptype{none} % do not increment counter
\begin{longtable}[]{@{}
  >{\raggedright\arraybackslash}p{(\linewidth - 4\tabcolsep) * \real{0.3333}}
  >{\raggedright\arraybackslash}p{(\linewidth - 4\tabcolsep) * \real{0.3333}}
  >{\raggedright\arraybackslash}p{(\linewidth - 4\tabcolsep) * \real{0.3333}}@{}}
\toprule\noalign{}
\begin{minipage}[b]{\linewidth}\raggedright
Indikator Evaluasi
\end{minipage} & \begin{minipage}[b]{\linewidth}\raggedright
Nilai
\end{minipage} & \begin{minipage}[b]{\linewidth}\raggedright
Keterangan
\end{minipage} \\
\midrule\noalign{}
\endhead
\bottomrule\noalign{}
\endlastfoot
\textbf{Total Kasus Tercatat} & \textbf{32 Kasus} & Dinamika Konflik di
Sulawesi (KPA/LSM) \\
\end{longtable}
}

\begin{figure}
\centering
\pandocbounded{\includegraphics[keepaspectratio,alt={Timeline Letusan Konflik Agraria}]{visuals_bab7/chart_7_2_konflik_hukum.png}}
\caption{Timeline Letusan Konflik Agraria}
\end{figure}

\begin{center}\rule{0.5\linewidth}{0.5pt}\end{center}

\subsection{7.3 Inkonsistensi Iklim: Karpet Merah PLTU
Captive}\label{inkonsistensi-iklim-karpet-merah-pltu-captive}

\textbf{Paradoks Hilirisasi Hijau dan Karpet Merah untuk PLTU Batubara
Captive}

\textbf{Metode: Penyaringan Agregat Dataset Eksternal (Global Coal Plant
Tracker GEM)}

\subsubsection{Metodologi: Agregasi Beban Karbon PLTU
Captive}\label{metodologi-agregasi-beban-karbon-pltu-captive}

\textbf{Metode Analisis:} Sub-bab ini menggunakan inventarisasi agregat
kuantitatif (\emph{Quantitative Inventory Aggregation}) dari database
global PLTU batubara independen (\emph{captive}).

\begin{enumerate}
\def\labelenumi{\arabic{enumi}.}
\tightlist
\item
  \textbf{Model Ekstraksi Kapasitas Fosil:}

  \begin{itemize}
  \tightlist
  \item
    \textbf{Isolasi Regional:} Melakukan pemfilteran data inventaris
    energi kotor (PLTU) yang berlokasi secara presisi di kawasan
    industri strategis pulau Sulawesi.
  \item
    \textbf{Kuantifikasi Kontradiksi Karbon:} Menghitung total jumlah
    \emph{unit} pembangkit dan agregat luaran listrik kotor (dalam
    satuan Megawatt) yang dibangun secara masif demi menopang pabrik
    pemurnian nikel, yang notabene dipromosikan sebagai proyek energi
    ramah lingkungan.
  \end{itemize}
\item
  \textbf{Kalkulasi/Formula Pengolahan:}

  \begin{itemize}
  \tightlist
  \item
    \texttt{Total\_Beban\_Karbon\ =\ SUM(Capacity\_MW)\ GROUP\ BY\ Provinsi}
  \item
    \texttt{Total\_Infrastruktur\_Kotor\ =\ COUNT(Unit\_PLTU)}
  \end{itemize}
\item
  \textbf{Variabel \& Fitur Data:}

  \begin{itemize}
  \tightlist
  \item
    \textbf{Spesifikasi Pembangkit:} \texttt{Capacity\ (MW)},
    \texttt{Start\ year}, \texttt{Provinsi\ (Subnational\ unit)}
  \end{itemize}
\item
  \textbf{Dataset \& File:}

  \begin{itemize}
  \tightlist
  \item
    \texttt{data/processed/sulawesi\_pltu\_captive.csv}
  \end{itemize}
\end{enumerate}

\begin{center}\rule{0.5\linewidth}{0.5pt}\end{center}

Komitmen transisi energi global dan pengembangan rantai pasok industri
nikel memegang peranan strategis. Di saat yang sama, pemenuhan kebutuhan
energi untuk fasilitas pengolahan nikel (\emph{smelter}) di Sulawesi
masih didominasi oleh Pembangkit Listrik Tenaga Uap (PLTU) Batubara
\emph{Captive}.

Data dari \emph{Global Coal Plant Tracker} (GEM) mencatat keberadaan
unit PLTU \emph{Captive} yang beroperasi maupun direncanakan di kawasan
industri Sulawesi Tengah dan Sulawesi Tenggara. Pemanfaatan energi
berbasis batu bara pada industri ini menghasilkan tantangan tersendiri
bagi pengelolaan emisi gas rumah kaca dan kualitas udara ambien.

Kondisi ini menunjukkan perlunya strategi percepatan transisi energi
bersih di sektor industri ekstraktif guna menyelaraskan target
hilirisasi dengan komitmen penurunan emisi nasional.

{\def\LTcaptype{none} % do not increment counter
\begin{longtable}[]{@{}
  >{\raggedright\arraybackslash}p{(\linewidth - 4\tabcolsep) * \real{0.3333}}
  >{\raggedright\arraybackslash}p{(\linewidth - 4\tabcolsep) * \real{0.3333}}
  >{\raggedright\arraybackslash}p{(\linewidth - 4\tabcolsep) * \real{0.3333}}@{}}
\toprule\noalign{}
\begin{minipage}[b]{\linewidth}\raggedright
Indikator PLTU Captive
\end{minipage} & \begin{minipage}[b]{\linewidth}\raggedright
Nilai
\end{minipage} & \begin{minipage}[b]{\linewidth}\raggedright
Keterangan
\end{minipage} \\
\midrule\noalign{}
\endhead
\bottomrule\noalign{}
\endlastfoot
\textbf{Total Unit PLTU Captive} & \textbf{67 Unit} &
Beroperasi/Dibangun/Direncanakan Khusus Kawasan Sulawesi \\
\textbf{Total Kapasitas Pembangkitan Kotor} & \textbf{12,245 MW} &
Pembangkit Fosil Penopang Smelter \\
\end{longtable}
}

\begin{figure}
\centering
\pandocbounded{\includegraphics[keepaspectratio,alt={Timeline Pertumbuhan Kapasitas PLTU Captive}]{visuals_bab7/chart_7_3_pltu_captive_timeline.png}}
\caption{Timeline Pertumbuhan Kapasitas PLTU Captive}
\end{figure}

\end{document}
